\documentclass[10pt,a4paper]{amsart}
\usepackage{amsmath,amssymb,amsthm,graphicx}
% For commutative diagrams we use Xy-pic
\usepackage[all]{xy}

% We use 2cell for 2-commutative diagrams.
\xyoption{2cell}
% Theorem environments.
%
\theoremstyle{plain}
\newtheorem{theorem}[subsection]{Theorem}
\newtheorem{proposition}[subsection]{Proposition}
\newtheorem{lemma}[subsection]{Lemma}

\theoremstyle{definition}
\newtheorem{definition}[subsection]{Definition}
\newtheorem{example}[subsection]{Example}
\newtheorem{exercise}[subsection]{Exercise}
\newtheorem{situation}[subsection]{Situation}

\theoremstyle{remark}
\newtheorem{remark}[subsection]{Remark}
\newtheorem{remarks}[subsection]{Remarks}

\numberwithin{equation}{subsection}

% Macros
%
\def\lim{\mathop{\mathrm{lim}}\nolimits}
\def\colim{\mathop{\mathrm{colim}}\nolimits}
\def\Spec{\mathop{\mathrm{Spec}}}
\def\Hom{\mathop{\mathrm{Hom}}\nolimits}
\def\Ext{\mathop{\mathrm{Ext}}\nolimits}
\def\SheafHom{\mathop{\mathcal{H}\!\mathit{om}}\nolimits}
\def\SheafExt{\mathop{\mathcal{E}\!\mathit{xt}}\nolimits}
\def\Sch{\mathit{Sch}}
\def\Mor{\mathop{\mathrm{Mor}}\nolimits}
\def\Ob{\mathop{\mathrm{Ob}}\nolimits}
\def\Sh{\mathop{\mathit{Sh}}\nolimits}
\def\NL{\mathop{N\!L}\nolimits}
\def\CH{\mathop{\mathrm{CH}}\nolimits}
\def\proetale{{pro\text{-}\acute{e}tale}}
\def\etale{{\acute{e}tale}}
\def\QCoh{\mathit{QCoh}}
\def\Ker{\mathop{\mathrm{Ker}}}
\def\Im{\mathop{\mathrm{Im}}}
\def\Coker{\mathop{\mathrm{Coker}}}
\def\Coim{\mathop{\mathrm{Coim}}}

% Boxtimes
%
\DeclareMathSymbol{\boxtimes}{\mathbin}{AMSa}{"02}

%
% Macros for moduli stacks/spaces
%
\def\QCohstack{\mathcal{QC}\!\mathit{oh}}
\def\Cohstack{\mathcal{C}\!\mathit{oh}}
\def\Spacesstack{\mathcal{S}\!\mathit{paces}}
\def\Quotfunctor{\mathrm{Quot}}
\def\Hilbfunctor{\mathrm{Hilb}}
\def\Curvesstack{\mathcal{C}\!\mathit{urves}}
\def\Polarizedstack{\mathcal{P}\!\mathit{olarized}}
\def\Complexesstack{\mathcal{C}\!\mathit{omplexes}}
% \Pic is the operator that assigns to X its picard group, usage \Pic(X)
% \Picardstack_{X/B} denotes the Picard stack of X over B
% \Picardfunctor_{X/B} denotes the Picard functor of X over B
\def\Pic{\mathop{\mathrm{Pic}}\nolimits}
\def\Picardstack{\mathcal{P}\!\mathit{ic}}
\def\Picardfunctor{\mathrm{Pic}}
\def\Deformationcategory{\mathcal{D}\!\mathit{ef}}
\usepackage{cleveref}
\crefname{originallemma}{original lemma}{original lemmas}
\Crefname{originallemma}{Original lemma}{Original lemmas}
\makeatletter
\newlabel{lemma-locally-constant-on-connected}{{0DV5}{source}{}{sourcepointer-1}{}}
\newlabel{lemma-locally-constant-on-connected@cref}{{[originallemma][0][]0DV5}{[1][0][]source}{}{}{}}
\makeatother
\crefname{originallemma}{original lemma}{original lemmas}
\Crefname{originallemma}{Original lemma}{Original lemmas}
\makeatletter
\newlabel{pione-section-finite-etale-under-galois}{{03SF}{source}{}{sourcepointer-2}{}}
\newlabel{pione-section-finite-etale-under-galois@cref}{{[originallemma][0][]03SF}{[1][0][]source}{}{}{}}
\makeatother
\crefname{originalremark}{original remark}{original remarks}
\Crefname{originalremark}{Original remark}{Original remarks}
\makeatletter
\newlabel{remark-functorial-locally-constant-on-connected}{{0DV6}{source}{}{sourcepointer-3}{}}
\newlabel{remark-functorial-locally-constant-on-connected@cref}{{[originalremark][0][]0DV6}{[1][0][]source}{}{}{}}
\makeatother
\title{Making a locally constant F-l sheaf constant after a finite etale cover of degree prime to l}
\author{The Stacks Project Authors}
\date{stacks-project:etale-cohomology.tex:local-snapshot-20261011; GNU FDL 1.2 or later}
\begin{document}
\maketitle
\noindent\textit{Editorial note.} All delivered mathematical text is the author's own, copied from the pinned Stacks Project chapter file etale-cohomology.tex; no statement or proof step was rewritten, reordered or removed.\par\medskip
\noindent\textit{Editorial note.} The proof is the author's complete argument; the covering, its degree and the filtration are all produced inside the counted span.\par\medskip
\noindent\textit{Editorial note.} Every labelled reference occurring inside the counted spans is carried in assembly.external\_references as a non-load-bearing navigation pointer (same-chapter ones at their definition site, cross-chapter ones at the invocation inside the counted span), so no load-bearing mathematics is replaced by a stub. No bibliography citation occurs inside the counted spans.\par\medskip
\noindent\textit{Editorial note.} The checked chapter list for same-name collisions is the set of stacks chapters staged in this repository: chow, cohomology, divisors, etale-cohomology, flat, more-algebra, more-morphisms, morphisms, spaces-more-morphisms, varieties together with this card's chapter etale-cohomology.\par\medskip
\noindent\textit{Editorial note.} The preamble is the chapter's own native preamble excerpt already used for this source: the Xy-pic setup (native preamble lines 14-18) and the native theorem declarations with the mathematical macros (native preamble lines 154-217).\par\medskip
\noindent\textit{Allowed standard prerequisites (editorial).}\begin{itemize}
\item Finite type locally constant sheaves of F\_l-vector spaces and the associated finite representations of the fundamental group. Same chapter, the classification of finite type locally constant sheaves of F\_l-vector spaces by continuous representations of the etale fundamental group on finite dimensional F\_l-vector spaces. Supplies the translation of the problem into a statement about a finite group representation.
\item Existence of finite etale covers trivialising a finite locally constant sheaf. Same chapter, the description of a locally constant sheaf as constant over the members of a finite etale covering. Supplies the covering over which the sheaf becomes constant.
\item Representations of a finite group of order prime to l and their lifts. Same chapter together with the group-theoretic input used by the author: over a finite group of order prime to l a finite dimensional representation in characteristic l is entirely decomposable into characters, so after a suitable cover the filtration by trivial quotients becomes visible. Supplies the argument producing the filtration with quotients Z/l.
\end{itemize}\medskip
\ifdefined\hypertarget\hypertarget{sourcepointer-1}{}\fi
\noindent\textit{Source reference.} Original lemma 0DV5: sources/stacks-etale-cohomology/pinned-original-source.tex, source lines 10118--10118. Navigation pointer invoked at proof line 10426; the counted span contains the pointer only and the cited result is not reproduced inside it, so the pointer is registered as a non-load-bearing navigation stub in the established form for Stacks material. No bibliography citation is involved.\par\medskip
\ifdefined\hypertarget\hypertarget{sourcepointer-2}{}\fi
\noindent\textit{Source reference.} Original lemma 03SF: sources/stacks-etale-cohomology/pinned-original-source.tex, source lines 10437--10437. Navigation pointer invoked at proof line 10437; the counted span contains the pointer only and the cited result is not reproduced inside it, so the pointer is registered as a non-load-bearing navigation stub in the established form for Stacks material. No bibliography citation is involved.\par\medskip
\ifdefined\hypertarget\hypertarget{sourcepointer-3}{}\fi
\noindent\textit{Source reference.} Original remark 0DV6: sources/stacks-etale-cohomology/pinned-original-source.tex, source lines 10294--10294. Navigation pointer invoked at proof line 10446; the counted span contains the pointer only and the cited result is not reproduced inside it, so the pointer is registered as a non-load-bearing navigation stub in the established form for Stacks material. No bibliography citation is involved.\par\medskip
\begin{lemma}
\label{lemma-pullback-filtered}
Let $S$ be a connected scheme. Let $\ell$ be a prime number. Let
$\mathcal{F}$ be a finite type, locally constant sheaf of
$\mathbf{F}_\ell$-vector spaces on $S_\etale$.
Then there exists a finite \'etale morphism
$f : T \to S$ of degree prime to $\ell$ such that $f^{-1}\mathcal{F}$
has a finite filtration whose successive quotients are
$\underline{\mathbf{Z}/\ell\mathbf{Z}}_T$.
\end{lemma}

\begin{proof}
Choose a geometric point $\overline{s}$ of $S$.
Via the equivalence of Lemma \ref{lemma-locally-constant-on-connected}
the sheaf $\mathcal{F}$ corresponds to a finite dimensional
$\mathbf{F}_\ell$-vector space $V$ with a continuous
$\pi_1(S, \overline{s})$-action.
Let $G \subset \text{Aut}(V)$ be the image of the homomorphism
$\rho : \pi_1(S, \overline{s}) \to \text{Aut}(V)$ giving the action.
Observe that $G$ is finite.
The surjective continuous homomorphism
$\overline{\rho} : \pi_1(S, \overline{s}) \to G$
corresponds to a Galois object $Y \to S$ of
$\textit{F\'Et}_S$ with automorphism group $G = \text{Aut}(Y/S)$, see
Fundamental Groups, Section \ref{pione-section-finite-etale-under-galois}.
Let $H \subset G$ be an $\ell$-Sylow subgroup.
We claim that $T = Y/H \to S$ works. Namely, let $\overline{t} \in T$
be a geometric point over $\overline{s}$. The image of
$\pi_1(T, \overline{t}) \to \pi_1(S, \overline{s})$
is $(\overline{\rho})^{-1}(H)$ as follows from the functorial
nature of fundamental groups. Hence the action of $\pi_1(T, \overline{t})$
on $V$ corresponding to $f^{-1}\mathcal{F}$ is through
the map $\pi_1(T, \overline{t}) \to H$, see
Remark \ref{remark-functorial-locally-constant-on-connected}. As
$H$ is a finite $\ell$-group, the irreducible constituents of the
representation $\rho|_{\pi_1(T, \overline{t})}$
are each trivial of rank $1$ (this is a simple lemma on
representation theory of finite groups; insert future reference here).
Via the equivalence of
Lemma \ref{lemma-locally-constant-on-connected}
this means $f^{-1}\mathcal{F}$ is a successive extension of
constant sheaves with value $\underline{\mathbf{Z}/\ell\mathbf{Z}}_T$.
Moreover the degree of $T = Y/H \to S$ is prime to $\ell$
as it is equal to the index of $H$ in $G$.
\end{proof}

\end{document}
